% SOURCE docs/release/2026-10-29/source/docs/writing-repair/2026-10-29/src/8066f27059ec873c.tex CHAR 0
\chapter{Time series, structural shocks and observation frequency}
\label{r3:time-series}
\section{A temporal prediction is not yet a shock intervention}
In treasury and deposit management, a bank may ask how a policy-rate
shock changes balances over the next year. In collections it may ask
how repeated assistance changes repayment along a customer trajectory.
Both have time indices, but the interventions differ. A structural
macroeconomic shock perturbs one innovation while the system's other
mechanisms continue to respond. A sustained policy replaces a sequence
of decision rules. A forecast conditions on observed historical values.
These are different operations on a dynamic model.

Let $X_t$ be a vector of observed series and let $\mathcal F_t$ be
their observed past. Granger noncausality of series $X^j$ for $X^i$
is a restriction saying that removing the past of $X^j$ from the
conditioning information does not change the conditional law, or in
its mean version the conditional mean, of $X^i_t$. It is a property
of a specified information set. It acquires an intervention
interpretation only under extra conditions on hidden causes, timing,
instantaneous effects and structural innovations
\citep[Chapter 10, Section 10.3]{peters2017}.

For a simple counterexample, let independent $U_t$ cause
$X_t=U_t+\eta_t$ and $Y_{t+1}=U_t+\nu_{t+1}$, with independent
noise. The past of $X$ helps predict $Y$ even though intervening on
the measurement $X_t$ leaves $U_t$ and $Y_{t+1}$ unchanged.
Conversely, an entirely instantaneous effect of $X_t$ on $Y_t$ with
independent observations over time need not generate any predictive
benefit from lagged $X$. Neither failure is repaired by a more
accurate nonlinear prediction model.

\section{Reduced and structural dynamics contain different information}
\label{r3:svar}
Consider a stable reduced-form VAR(1),
\[
 X_t=FX_{t-1}+u_t,\qquad \E[u_tu_t']=\Sigma_u.
\]
Suppose a structural model writes $u_t=B\varepsilon_t$ with
independent, unit-variance structural shocks. Under causal autonomy,
the response at horizon $h$ to a unit shock in coordinate $j$ is
\begin{equation}
 \operatorname{IRF}_j(h)=F^hBe_j.
 \label{r3:structural-irf}
\end{equation}
Iterate the recursion after changing only $\varepsilon_{j,t}$ to
prove this formula. It presumes that $F$ and the other mechanisms
remain invariant. Imposing a new interest-rate rule changes a mechanism
and can change $F$; that policy counterfactual is not generally a
single historical impulse response.

Gaussian likelihood identifies $F$ and $BB'$, but for any orthogonal
$Q$, $BQ(BQ)'=BB'$. Independent Gaussian shocks remain independent
after rotation. Hence identical observational dynamics permit
different shock responses unless restrictions attach an economic
meaning to a particular column. For example, with
$F=\operatorname{diag}(0.7,0.4)$ and $\Sigma_u=I$, both $B=I$
and
\[
 B=\frac1{\sqrt2}\begin{pmatrix}1&-1\\1&1\end{pmatrix}
\]
fit the same Gaussian process. The first shock affects only the first
series in one representation and both in the other. Normalizing its
impact on the first series to one gives horizon-$h$ second-series
responses $0$ and $0.4^h$. Even knowing the entire observed
distribution cannot choose that financial interpretation.

Recursive zero restrictions, externally justified instruments,
heteroskedastic regimes, sign restrictions and non-Gaussian independent
innovations use different information to restrict $B$. Sign restrictions
can leave a set of admissible rotations, so the target can remain set
identified. A prior over those rotations gives a model-dependent
distribution over the set; it does not turn the restrictions into a
point-identification theorem.

\subsection{Non-Gaussianity adds information, with substantive costs}
For the model
$X_t=B_0X_t+\sum_{\ell=1}^pB_\ell X_{t-\ell}+\varepsilon_t$,
reduced lag matrices are $(I-B_0)^{-1}B_\ell$.
\citet[Sections 2 and 5]{hyvarinen2010} uses mutually independent,
non-Gaussian innovations over components and time, and an acyclic
instantaneous graph, to identify and estimate the structural system.
ICA identifies the innovation mixing up to its scale/permutation
ambiguities, and the acyclic structural restrictions resolve the
appropriate normalization and ordering. The assumptions add higher-order
information beyond covariance. They also exclude latent common causes
in that model. Fat-tailed bank returns alone do not establish
independent structural innovations, and stochastic volatility can
produce marginal non-Gaussianity with shared dependence.

The distinction between $B_\ell$ and $(I-B_0)^{-1}B_\ell$ explains
why a reduced lag coefficient includes instantaneous propagation.
Reporting it as one direct structural edge can misstate both its
size and its mechanism. Non-Gaussian discovery is therefore a
scientific identification route with testable implications and
untestable restrictions, not a preprocessing option for an otherwise
unidentified VAR.

\section{External instruments identify a normalized dynamic response}
\label{r3:lpiv}
Let a structural moving-average representation be
\[
 X_t=\sum_{j=0}^\infty\Theta_j\varepsilon_{t-j},
\]
with square-summable coefficients and orthogonal zero-mean shocks.
Suppose a scalar external instrument $Z_t$ obeys
\begin{align*}
 \E[Z_t\varepsilon_{1t}]&=\alpha\ne0,\\
 \E[Z_t\varepsilon_{kt}]&=0\quad(k\ne1),\\
 \E[Z_t\varepsilon_{k,t+j}]&=0\quad(j\ne0,\ \text{all }k).
\end{align*}
The last exclusion concerns both leads and lags. Since $X_{t+h}$
contains many past shocks, contemporaneous instrument validity alone
does not justify a direct horizon regression. Multiplying the
moving-average representation by $Z_t$ leaves only the date-$t$
target shock, giving
\begin{equation}
 \frac{\E[Z_tX_{i,t+h}]}{\E[Z_tX_{1,t}]}
       =\frac{\Theta_{h,i1}}{\Theta_{0,11}}.
 \label{r3:lpiv-ratio}
\end{equation}
With unit-impact normalization $\Theta_{0,11}=1$, the ratio is the
structural response in the units of the policy indicator.
This reconstructs Condition LP-IV and equation (8) of
\citet[Sections 2.2--2.3]{stock2018}.

If the instrument contains lagged structural shocks, controls must
actually remove the corresponding covariance while preserving
relevance. Merely adding several lags because a package defaults to
them does not prove the residualized exclusions. A monetary surprise
that also reveals central-bank information about future fundamentals
may load on another shock; a narrow measurement window alone does
not settle that exclusion. The target and normalization also matter:
for a news shock with zero immediate policy-indicator response,
the denominator at $h=0$ is zero. A justified future-impact
normalization defines a different scale
\citep[Sections 2.4--2.5]{stock2018}.

\subsection{Direct projections and VAR propagation}
An LP-IV estimates a separate instrumented regression for each
horizon. A proxy-SVAR first recovers reduced innovations and dynamics,
then identifies a structural impact column from instrument--innovation
covariances and propagates it. Under suitable invertibility,
structural shocks are recoverable from present and past observed
innovations. That condition is additional information about the
observation system. The direct LP-IV argument above can identify a
response without a finite-order invertible VAR, but its lead/lag
exclusion is strong. Under different control-variable conditions
the requirements of the two approaches can coincide
\citep[Sections 2--3]{stock2018}.

A scalar illustration shows what noninvertibility loses. Let
$X_t=\varepsilon_t+\theta\varepsilon_{t-1}$ with Gaussian white
noise of variance $\sigma^2$ and $|\theta|>1$. Its autocovariances
are identical to those of the invertible representation with
coefficient $1/\theta$ and innovation variance $\theta^2\sigma^2$:
both have variance $(1+\theta^2)\sigma^2$ and lag-one covariance
$\theta\sigma^2$. Higher-lag covariances vanish, so the Gaussian
laws agree. Their unit-impact structural responses at horizon one
are nevertheless $\theta$ and $1/\theta$. Choosing the invertible
likelihood representation is a normalization of observable prediction;
calling its innovation the economic shock needs justification.

LP estimators have more freedom across horizons; VAR propagation
imposes restrictions through powers of common lag matrices. This
creates an estimation bias/variance comparison after identification,
not a theorem that direct projections always answer a causal question
or that VARs always have lower error. The constructions and inference
alternatives are reviewed by \citet[Sections 2 and 6]{jorda2024}.
Moreover, identifying the normalized response does not automatically
identify the shock's variance, every historical shock realization,
or its share of forecast-error variance. Those require additional
recoverability and normalization information \citep[Section 2.5]{stock2018}.

\section{Observation frequency is part of the model}
\label{r3:temporal-observation}
For the latent or fully observed linear state
$S_{t+1}=FS_t+\eta_{t+1}$ with independent innovations of covariance
$Q$, observing only every $m$th state gives
\begin{equation}
 S_{t+m}=F^mS_t+\sum_{j=0}^{m-1}F^j\eta_{t+m-j},\qquad
 Q_m=\sum_{j=0}^{m-1}F^jQ(F^j)'.
 \label{r3:subsample}
\end{equation}
Repeated substitution proves both formulas. The coarse observations
identify the pair $(F^m,Q_m)$ under their statistical model; recovering
$F,Q$ is a new inverse problem. For a scalar example, daily
coefficients $0.7$ and $-0.7$ with equal innovation variance produce
the same stationary Gaussian law when sampled every second day:
both have coarse coefficient $0.49$ and noise variance $1.49Q$.
Their one-day intervention responses have opposite signs. Additional
fine-frequency data or structural restrictions are needed to recover
that difference.

Aggregation is different from subsampling. If
$x_t=\rho x_{t-1}+\epsilon_t$ and a reported flow is
$s_k=x_{2k-1}+x_{2k}$, direct algebra gives
\begin{equation}
 s_k=\rho^2s_{k-1}+v_k,\qquad
 v_k=\rho\epsilon_{2k-2}+(1+\rho)\epsilon_{2k-1}+\epsilon_{2k}.
 \label{r3:aggregate}
\end{equation}
Adjacent $v_k$ share an innovation, so
$\Cov(v_k,v_{k-1})=\rho\sigma_\epsilon^2$ and
$v_k$ also correlates with $s_{k-1}$. Fitting an ordinary AR(1)
to the aggregates does not recover the fine-scale coefficient by
taking a square root. Monthly average balances, period-end stocks
and summed cash flows require different observation operators.
Temporal aggregation can therefore change conditional independences
used by a discovery procedure, even when no economic mechanism changes.

In a linear continuous-time controlled model
$dS_t=(FS_t+Ba_t)dt+G\,dW_t$, the integrating-factor solution over
an interval $\Delta$ with constant action is
\begin{align}
 S_{t+\Delta}&=e^{F\Delta}S_t+
 \left(\int_0^\Delta e^{Fs}B\,ds\right)a_t+\eta_\Delta,\nonumber\\
 \Var(\eta_\Delta)&=\int_0^\Delta e^{Fs}GG'e^{F's}\,ds.
 \label{r3:continuous-discrete}
\end{align}
Thus irregular intervals change transition, action and noise kernels.
They cannot be handled merely by assigning consecutive integer
indices to timestamps. Different matrix logarithms or aliased
oscillations can also complicate recovery of $F$ from discrete
transitions. The derivation is a linear-model calculation; it does
not establish a general continuous-time causal identification theorem.

When only emissions $O_t=CS_t+\nu_t$ are observed, an invertible
coordinate change $S_t^*=RS_t$ produces transformed kernels
$F^*=RFR^{-1}$, $B^*=RB$, $C^*=CR^{-1}$ and $Q^*=RQR'$.
The observation law and coherently transformed input/output responses
are unchanged. This state-coordinate ambiguity is harmless for those
responses. It differs from the subsampling or shock ambiguity above,
where the requested intervention response changes. An estimator must
distinguish representational nonuniqueness from causal nonidentification.

\section{Dependent inference must follow the score process}
\label{r3:series-inference}
For a fixed horizon $h$, let $m_{t,h}(\beta_h)$ be an identified
moment and $J_h=\E[\partial_\beta m_{t,h}]$ nonsingular.
Under a stationary score CLT and a negligible estimation remainder,
\[
 \sqrt T(\widehat\beta_h-\beta_h)
 =-J_h^{-1}T^{-1/2}\sum_t m_{t,h}(\beta_h)+o_p(1).
\]
The long-run score covariance is
\begin{equation}
 S_h=\sum_{k=-\infty}^{\infty}
             \Cov(m_{t,h},m_{t-k,h}),\qquad
 \operatorname{avar}(\sqrt T\widehat\beta_h)
                  =J_h^{-1}S_hJ_h^{-\prime}.
 \label{r3:hac}
\end{equation}
Here the score CLT, covariance summability and a nondegenerate variance
are imposed conditions; stationarity alone does not imply them.
Covariance estimation needs further conditions;
a kernel HAC estimator truncates weighted empirical lag covariances
at a bandwidth that grows suitably slowly with $T$. Specifying that
estimator without stating its regime is not an uncertainty guarantee.

Multi-step regression errors often overlap in future innovations, but
it is the instrument-times-error score that controls inference. Under
additional orthogonality, a lag-augmented construction can yield
serially uncorrelated scores even with overlapping regression errors.
The ordinary and augmented constructions therefore need different
justifications; \citet[Sections 6.1--6.3]{jorda2024} discusses these
distinctions and the source-specific uniform results. An arbitrary
nonlinear LP with a learned instrument does not inherit them.

Fitting nuisance functions on separate time blocks likewise requires
a dependence argument. Under absolute regularity, coupling a block
to an independent version can have an error controlled by the
absolute-regularity coefficient at the separating gap. With many
blocks, the accumulated error must vanish. The gap must also account
for an outcome that extends $h$ dates after its origin; disjoint
origin dates can still share outcomes and innovations. Random row
splitting supplies neither control. A prospective evaluation additionally
requires training and features available at the historical decision
date; an offline stationary-law cross-fitting argument alone does
not certify that prospective use.

Joint inference across horizons needs the stacked score covariance,
including cross-horizon terms. A collection of pointwise 95\% intervals
is not a simultaneous 95\% response band. Weak normalization denominators,
unit roots, structural breaks, growing horizons and heavy tails can
each change the limiting argument. The appropriate result may be an
identified set, a sensitivity analysis or a shorter supported horizon
rather than a smooth but unjustified response curve.



% SOURCE docs/release/2026-10-29/source/docs/writing-repair/2026-10-29/src/6b090a7a6eb36c17.tex CHAR 0
\section{A complete stationary response experiment}
\label{s4:response-construction}
The moment identity in Equation~\eqref{r3:lpiv-ratio} becomes a usable
operation only after specifying how observations repeat. Consider a
three-dimensional vector containing a policy rate, deposits and lending,
all in declared units. Let
\[
 X_t=FX_{t-1}+B\varepsilon_t,\qquad
 Z_t=a\varepsilon_{1t}+v_t,\qquad
 v_t=\rho_vv_{t-1}+\sqrt{1-\rho_v^2}\zeta_t.
\]
The innovations are mutually independent standard Gaussian vectors/scalars;
the joint process starts in its stationary law. The eigenvalues of $F$ and
$\rho_v$ have modulus below one, with a fixed stability margin. The matrix
$B$ is nonsingular, $B_{11}\ne0$, and initially $aB_{11}$ is bounded away
from zero. Gaussianity is a statistical restriction of this reference
experiment, while the instrument restrictions supply causal identification.
The Gaussian observational VAR itself does not identify $B$.

Fix $H$ before sampling, and use origins $t=1,\ldots,n=T-H$ for every
coordinate. Index $j=(i,h)$ for $i\in\{2,3\}$ and $0\le h\le H$.
Set $U_{tj}=Z_tX_{i,t+h}$, $D_t=Z_tX_{1t}$,
$q=\E D_t=aB_{11}$, and $\theta_j=\E U_{tj}/q$.
The estimate and its exact linearization are
\begin{equation}
 \widehat\theta_j=\bar U_j/\bar D,\qquad
 \sqrt n(\widehat\theta_j-\theta_j)
 =\frac{1}{\bar D}\frac1{\sqrt n}\sum_{t=1}^n
       Z_t(X_{i,t+h}-\theta_jX_{1t}).
 \label{s4:ratio-linearization}
\end{equation}
This identity avoids a hidden regression remainder. A one-percentage-point
impact normalization is meaningful only if the policy rate is expressed
in percentage points. A change to basis points rescales the response.

The score $M_{tj}=U_{tj}-\theta_jD_t$ has mean zero. The long-run matrix
$\Omega=\sum_{k\in\mathbb Z}\Cov(M_t,M_{t-k})$ contains both serial
and cross-horizon dependence. In particular, independence of $v$ and $X$
does not make their product serially independent. Its contribution to
the lag covariance is
$\Cov(v_t,v_{t-k})\Cov(X_{i,t+h}-\theta_jX_{1t},
X_{i',t-k+h'}-\theta_{j'}X_{1,t-k})$.
Persistent measurement noise in an otherwise valid instrument can therefore
matter greatly for standard errors. This is the score-level distinction
emphasized by \citet[Section 2.3]{stock2018}.

\subsection{Why a central limit theorem applies here}
For this particular experiment the limiting argument can be reconstructed
directly. The jointly Gaussian vector process $(X_t,Z_t)$ has geometrically
decaying autocovariances and a bounded spectral density. Gaussian moment
pairing shows that the covariance of any two centered quadratic products
is a sum of products of two covariances. Hence the score covariances are
absolutely summable; a fixed horizon shift changes the constants, not
summability.

For any fixed vector $c$, collect the finitely many Gaussian observations
needed for $\sum_tc'M_t$ into $G_n$. That sum is a centered quadratic
form $G_n'A_nG_n-\operatorname{tr}(A_n\Sigma_n)$. Each observation is
paired with only a fixed number of lags, so $\|A_n\|$ is bounded.
Absolute covariance summability bounds $\|\Sigma_n\|$ uniformly. The
eigenvalues $\lambda_{nr}$ of
$\Sigma_n^{1/2}A_n\Sigma_n^{1/2}$ are consequently bounded, and the form
has the law $\sum_r\lambda_{nr}(\chi_{1r}^2-1)$ with independent terms.
Its variance divided by $n$ converges to $c'\Omega c$. If this limit is
positive, the fourth-moment Lyapunov ratio is bounded by
$C\sum_r\lambda_{nr}^4/(\sum_r\lambda_{nr}^2)^2=O(n^{-1})$.
If the limit is zero, the normalized form converges to zero in $L_2$.
Applying this argument to every $c$ proves the joint Gaussian limit.
With $\bar D\to q$, Equation~\eqref{s4:ratio-linearization} gives
$\sqrt n(\widehat\theta-\theta)\Rightarrow N(0,\Omega/q^2)$.
This is a project proof for a fixed-dimensional Gaussian linear law;
stationarity alone does not prove its conclusion.

\subsection{A covariance estimator with an explicit bandwidth}
Let $\widehat M_t=U_t-\widehat\theta D_t$ and center each score column.
With $L=\lfloor n^{1/4}\rfloor$, compute
\begin{align}
 \widehat\Gamma_k&=n^{-1}\sum_{t=k+1}^n
    (\widehat M_t-\bar M)(\widehat M_{t-k}-\bar M)',\nonumber\\
 \widehat\Omega_L&=\widehat\Gamma_0+
   \sum_{k=1}^L\left(1-\frac{k}{L+1}\right)
       (\widehat\Gamma_k+\widehat\Gamma_k'),\qquad
 \widehat V=\frac{\widehat\Omega_L}{n\bar D^2}.
 \label{s4:bartlett}
\end{align}
The common denominator $n$, including at nonzero lags, matters. Extending
the centered sequence by zeros, this estimator is the sum of outer
products of all length-$(L+1)$ moving sums, divided by $n(L+1)$. It is
positive semidefinite. The original construction and a general consistency
result are in \citet[equation (5), Theorems 1--2]{newey1986}; the latter's
general mixing theorem requires $L=o(n^{1/4})$. Our stated $n^{1/4}$
choice instead uses the stronger Gaussian structure just specified.

Here are the additional steps. Gaussian moment pairing bounds the variance
of every empirical lag covariance by $C/n$, uniformly in the included
lag: in each covariance of two fourth-degree terms, a surviving pairing
joins the two date groups, whose geometrically decaying covariance is
summable over their separation. There are finitely many pairings, and
shifts by $k$ do not change that sum. Summing the resulting standard
deviation bounds gives $O_p(L/\sqrt n)$. The Bartlett taper bias is
bounded by $(L+1)^{-1}\sum_k |k|\|\Gamma_k\|$ and the omitted tail
is geometric. Replacing the true parameter by a root-$n$ estimate changes
the sum by $O_p(L/\sqrt n)$, using Cauchy--Schwarz on each empirical
cross product. Centering and edge corrections vanish as well. Thus
\[
 \|\widehat\Omega_L-\Omega\|
 =O_p(L/\sqrt n)+O(L^{-1})+O(r^L),\qquad r<1,
\]
is a conservative bound in this fixed model. It proves consistency for
the stated bandwidth; it does not promise small bias at 800 monthly origins
or uniformly as $r$ approaches one.

For $d=2(H+1)$ response coordinates and positive marginal variances,
$\widehat\theta_j\pm z_{1-\alpha/(2d)}\sqrt{\widehat V_{jj}}$
is an asymptotic simultaneous Bonferroni band. A less conservative
alternative computes the $(1-\alpha)$ quantile of
$\max_j|G_j|$ for $G\sim N(0,\widehat R)$, where $\widehat R$ is
the correlation matrix of $\widehat V$. Numerical Gaussian draws
approximate that quantile; their Monte Carlo error is additional to the
sampling approximation. The experiments use 4,096 such draws and report
that limitation. A null-centered significance band tests a specified
zero-response hypothesis and should not be relabeled as a confidence
band for every response; see \citet[equations (5)--(8)]{inoue2023}.

\subsection{Weak instruments require a different result object}
Do not divide by a weak normalization to construct the inferential pivot.
For a proposed response $b$, test the zero mean of $U_{tj}-bD_t$.
Write $S_{UU},S_{UD},S_{DD}$ for the corresponding entries of the
Bartlett matrix of the centered joint products $(U_t,D_t)$. Accept $b$ if
\begin{equation}
 n(\bar U_j-b\bar D)^2
 \le c^2(S_{UU}-2bS_{UD}+b^2S_{DD}).
 \label{s4:weak-inversion}
\end{equation}
This is one quadratic inequality $A b^2+B b+C\le0$, with
$A=n\bar D^2-c^2S_{DD}$,
$B=-2n\bar U_j\bar D+2c^2S_{UD}$, and
$C=n\bar U_j^2-c^2S_{UU}$.
Its solution can be an interval, two unbounded rays, the whole line,
an empty set, or a degenerate singleton/half-line. These cases must survive
the interface. At the true $b$, the score CLT does not require dividing
by $q$, and remains applicable along bounded weak-instrument sequences
with a nondegenerate residual-score variance and fixed stability margin.
Bonferroni inversion across coordinates yields simultaneous coverage of
the true response vector. At $q=0$ the normalized causal response is
not identified by these moments; a statistical set cannot repair that.
No first-stage threshold is used to switch silently between guarantees.

\subsection{A fully specified propagation challenger}
Fit the VAR(1) by OLS, giving $\widehat F$, residuals $\widehat u_t$ and
$\widehat\kappa=n^{-1}\sum_tZ_t\widehat u_t$. Its response is
$e_i'\widehat F^h\widehat\kappa/\widehat\kappa_1$.
With $Q=\E X_{t-1}X_{t-1}'$, the influence matrix for $F$ is
$u_tX_{t-1}'Q^{-1}$. The influence vector for $\kappa$ is
$Z_tu_t-\kappa$ because the stipulated exclusions give
$\E[Z_tX_{t-1}]=0$; without that condition an additional derivative term
is required. For perturbations $(dF,d\kappa)$,
\[
 d\theta_{ih}=\frac{e_i'\sum_{r=0}^{h-1}F^r(dF)F^{h-1-r}\kappa}
                       {\kappa_1}
 +\frac{e_i'F^hd\kappa-\theta_{ih}d\kappa_1}{\kappa_1}.
\]
Apply this derivative to the joint score and use its joint HAC matrix.
This challenger propagates uncertainty in the fitted dynamics and impact
direction. Its dimensional savings rely on correct finite-order dynamics;
a misspecified VAR cannot borrow the direct moment estimator's robustness.

\section{Learned controls under explicit temporal separation}
\label{s4:dependent-learning}
Many bank analyses require hundreds of release-vintage controls. A
partially linear IV model supplies a precise learning target:
\[
 Y_t=\theta D_t+g(C_t)+e_t,\quad
 \E[e_t\mid C_t,Z_t]=0,\quad
 (\ell,r,m)(c)=\E[(Y_t,D_t,Z_t)\mid C_t=c].
\]
The orthogonal score and its identifying slope are
\[
 \psi_t(\theta,\eta)=(Z_t-m(C_t))
  \{Y_t-\ell(C_t)-\theta[D_t-r(C_t)]\},\qquad
 J=\E[(Z_t-m)(D_t-r)]\ne0.
\]
Instruments, controls and the outcome horizon must already satisfy the
causal restrictions. Good nuisance prediction does not supply them.
For fitted errors $\Delta m,\Delta\ell,\Delta r$, conditional on an
independent training sample, direct expansion gives the exact mean remainder
\begin{equation}
 \E\{\psi(\theta,\widehat\eta)-\psi(\theta,\eta)\}
 =\E[\Delta m(\Delta\ell-\theta\Delta r)].
 \label{s4:dependent-remainder}
\end{equation}
The first-order terms have zero means by their respective conditional
moment restrictions. This connects the dynamic procedure to the
orthogonality argument in Chapter~\ref{r3:learning} and
\citet{dml2018}, with a new sampling argument still needed.

One sufficient construction is deliberately explicit. Suppose the underlying
observables and fitted functions are bounded, the process is absolutely
regular with $\beta(k)\le C\exp(-ck)$, and lag/horizon windows are fixed.
Use two contiguous blocks of comparable size, separated by
$g=\lceil3\log n/c\rceil$ \emph{after} expanding both blocks for every
feature lag and future outcome. Train on either block and evaluate on the
other. Absolute regularity couples the pair of blocks to independent
blocks with error at most $C\exp(-cg)=O(n^{-3})$. Both fits use this
same coupled pair; independence of the two fitted estimates is not asserted.
An unknown $c$ is a scientific tuning uncertainty. The formula is a
sufficient regime specification, not an estimate of $c$ from a few scores.

Let $\delta_n$ bound the three population $L_2$ nuisance errors. For a
bounded centered score difference its lag covariance is bounded both by
$C\delta_n^2$ and by $C\exp(-ck)$. Splitting the lag sum at
$k\asymp\log(1/\delta_n)$ bounds its long-run variance by
$C\delta_n^2\{1+\log(1/\delta_n)\}$. Conditional centering, Chebyshev's
inequality and Equation~\eqref{s4:dependent-remainder} then give the oracle
linearization when
\[
 \delta_n\to0,\qquad
 \|\Delta m\|_2(\|\Delta\ell\|_2+|\theta|\|\Delta r\|_2)
       =o_p(n^{-1/2}).
\]
The oracle CLT can also be constructed here. Split the bounded centered
score into big blocks of length $m=\lfloor n^{1/3}\rfloor$, separated
by gaps $b=\lceil C\log n\rceil$ with $cC>2$. Summable score
covariances bound the variance of the discarded gaps by a constant times
their number of observations, $O(n^{2/3}\log n)=o(n)$. Coupling the
remaining big blocks to independent blocks fails with probability at most
$O(n^{2/3}e^{-cb})=o(1)$. Each centered block is bounded by $O(m)$,
which is $o(\sqrt n)$; hence the independent-block Lindeberg condition
holds. Their variance divided by $n$ converges to the summable long-run
variance. Cram\'er--Wold and the negligible gaps give the vector CLT.

For covariance estimation, a covariance of two score products decomposes
into a fourth cumulant plus two products of pairwise covariances. Bounded
exponential absolute regularity bounds the cumulant by an exponential in
the largest separating gap of the four dates; summing its possible gaps
is finite. Pairwise terms are summable as well. This gives a uniform
$C/n$ variance bound for empirical lag covariances, and the same
$O_p(L/\sqrt n)$ conservative stochastic bound as above. Across the deleted central gap,
do not pretend the two observed endpoints are adjacent. Sum lag products
within each evaluation block using the total retained length as denominator.
The missing oracle boundary pairs contribute $O(L^2/n)$, which vanishes
for $L=\lfloor n^{1/4}\rfloor$.

There is an additional condition for fitted-score covariance. If the
empirical root-mean-square score error is $r_n$, Cauchy--Schwarz bounds
the difference of the two HAC matrices by $O_p(Lr_n+Lr_n^2)$.
Take $\delta_n$ to be a deterministic high-probability bound on the
population errors. Conditional evaluation and Markov's inequality applied
to the mean squared score error give $r_n=O_p(\delta_n)$; the coupling
failure probability vanishes. Require $L\delta_n=o(1)$, as well as the product condition.
Taking all nuisance errors $o_p(n^{-1/4})$ is sufficient for the stated
bandwidth. A maximum lag-covariance error $o_p(1)$ alone cannot be
summed over an increasing number of lags.

\subsection{When sparsity supplies the required rates}
For a concrete learner suppose each conditional mean is a linear combination
of $p\le n^a$ bounded, standardized features, with at most $s$ nonzero
coefficients, plus an explicitly controlled approximation error. Assume
a population restricted eigenvalue bounded below. Partition training data
into alternating blocks of length $b=C_0\log n$, with
$C_0>(3+2a)/c$. Coupling alternate blocks, followed by bounded-block
concentration and a union bound over scores and Gram entries, gives
deviations of order $\lambda_n=C_1\sqrt{\log n\log(pn)/n}$ for a
constant $C_1$ depending on the envelope and dependence constants.
This conservative penalty includes a logarithmic dependence cost.

The Lasso basic inequality on the event
$\|n^{-1}\sum C_te_t\|_\infty\le\lambda_n/2$ restricts its error
to a cone. Restricted curvature then gives prediction error
$O_p(\sqrt{s}\lambda_n)$ and coefficient $\ell_1$ error
$O_p(s\lambda_n)$. Concentration of the Gram matrix transfers the
empirical prediction bound to population $L_2$ when $s\lambda_n=o(1)$.
Clipping predictions to the known outcome envelope cannot increase that
population error. Consequently
$s\log n\log(pn)=o(\sqrt n)$, with approximation errors obeying the
same product and bandwidth restrictions, supplies a complete sufficient
learning regime. This is substantially stronger than merely $p<n$.

\citet[Assumptions 3.1--3.4 and Theorem 3.1]{babii2020} develops a more
flexible sparse-group route under $\tau$-mixing and polynomial tails,
with a penalty depending explicitly on moments and decay rates. Its
prediction inequalities do not automatically supply the population nuisance
rates or the bounded coupling argument just used. A factor alternative
$C_t=\Lambda f_t+e_t^C$ economizes on a dense common component, but
dropping $e_t^C$ is valid adjustment only if the omitted component has no
remaining confounding role, or its bias is bounded. Sparsity and factor
reduction are therefore different scientific restrictions. Forecast RMSE
does not decide between them for causal inference.

The current proposal of \citet[Assumptions 2.1--2.5]{ciganovic2026}
uses reverse cross-fitting for time series and explicitly assumes
conditional plug-in stability. That assumption cannot be deduced from
disjoint row indices or reversibility. For example, a stationary Gaussian
process can be reversible and still have nonsummable long-memory
covariances. Its Supplement S1's increasing-bandwidth argument also needs
control of the accumulated fitted-score discrepancy. We retain the source
as a relevant 2026 alternative and use the stronger explicit separation
and rate construction above for the initial justified route. Removing
gaps requires a learner-specific stability proof.

\section{Common bank shocks determine the effective sample size}
\label{s4:date-panel}
For a fixed set of banks with target weights $w_i$, the relevant score at
date $t$ is $\bar\psi_t=\sum_iw_i\psi_{it}$, not an array of $NT$
independent rows. Apply the temporal construction to this aggregate when
its own conditions hold. For the transparent example
$\psi_{it}=C_t+e_{it}$, with independent idiosyncratic white noise,
\[
 \operatorname{avar}(\widehat\theta)
 =J^{-2}T^{-1}\left\{\Omega_C+\sigma_e^2\sum_iw_i^2\right\}.
\]
Equal weights reduce the idiosyncratic term by $N$ and leave the common
component unchanged. More banks do not create more monetary-policy shocks.
Sampling a random population of banks adds a separate uncertainty source;
time aggregation alone does not cover it. Staggered-treatment comparisons,
dynamic-panel moment restrictions and multiway-cluster methods from
Chapter~\ref{r3:panels} answer distinct questions and retain their own
assignment and asymptotic conditions.

\section{Persistence, cointegration and the level target}
\label{s4:integrated-responses}
An integrated balance series is economically plausible and falls outside
the stable-level experiment. For a VAR(1) with semisimple unit roots and
stable remaining roots,
\[
 \Delta X_t=(F-I)X_{t-1}+B\varepsilon_t
             =\alpha\beta'X_{t-1}+B\varepsilon_t.
\]
The rank of $F-I$ and the stationary error-correction combinations
$\beta'X_t$ describe long-run comovement. They do not identify which
innovation is a monetary, liquidity or credit-supply shock. The latter
still requires an instrument or another economic restriction.

There is an identified stationary transformation for the fixed-horizon
external-IV target. Suppose $\Delta X$ has a stationary Gaussian linear
representation with geometrically decaying coefficients and the original
lead/lag exclusions hold. Write $\theta_{-1}=0$. Then
\[
 \frac{\E[Z_t\Delta X_{i,t+h}]}{\E[Z_t\Delta X_{1t}]}
       =\theta_{ih}-\theta_{i,h-1},\qquad
 \theta_{ih}=\sum_{r=0}^h
       \frac{\E[Z_t\Delta X_{i,t+r}]}{\E[Z_t\Delta X_{1t}]}.
\]
The denominator still equals $aB_{11}$. The numerator sums to
$\E[Z_t(X_{i,t+h}-X_{i,t-1})]$: this is the response of the
\emph{level} to a one-time shock, recovered from stationary changes,
not the effect of permanently changing a rate rule. Estimate the stacked
change responses and premultiply by the block-diagonal cumulative-sum
matrix $C_H$. The level covariance is
$C_H\widehat V_\Delta C_H'$. For weak-IV inversion, first transform
the numerator moments and their full covariance, then invert; summing
separate confidence endpoints discards their dependence.

This route does not require estimating the cointegrating rank to define
the fixed-horizon functional. Efficient VECM fitting can use rank
restrictions as a challenger, with rank uncertainty retained. The state-space
connection is also substantive. In $T_{t+1}=T_t+\eta_{t+1}$,
$Y_{t+1}=BT_t+\epsilon_{t+1}$, the observed law identifies combinations
such as $B\Var(\eta)B'$, while invertible rotations of the latent trend
remain observationally equivalent. \citet[Theorems 1--2]{johansen2017}
shows that extracting the identified common-trend contribution and
estimating an individually normalized trend have different rate conditions.
Naming a filtered trend ``funding demand'' does not identify its intervention.

Lag augmentation offers a different persistence treatment. For
$Y_t=\phi Y_{t-1}+u_t$, regress $Y_{t+h}$ on $(Y_t,Y_{t-1})$.
The residualized first regressor approximates $u_t$ and the target
coefficient is $\phi^h$. Under mean independence of $u_t$ from the
innovations at \emph{all other dates}, the score formed with the
future projection error is serially uncorrelated even when that error
overlaps. \citet[Definition 1, Assumptions 1--3, Proposition 1]{lpaug2020}
proves marginal uniform inference for a specified VAR parameter space,
with eighth moments, summable innovation-product cumulants and scaled
design nonsingularity. Near unit roots its horizon condition is
$h_T/T\to0$. This theorem does not itself establish external-IV ratio
inference or simultaneous bands. We select the stationary-change IV route
for the integrated structural response and retain lag augmentation as
the justified alternative for its stated reduced-form projection target.

A regime break is a third problem. Pooling two different $F$ or impact
columns can average two economic mechanisms. A stationary transformation
does not remove a changing causal response. Segmented analyses require
enough dates in each segment, a declared treatment of estimated break
locations and, if used, a separate selection-aware uncertainty argument.
The product must identify that unresolved regime instead of interpreting
one pooled curve as a stable banking mechanism.


% SOURCE docs/release/2026-10-29/source/docs/writing-repair/2026-10-29/src/8066f27059ec873c.tex CHAR 14703


\section{Consequences for a large-scale time-series product}
The reusable object should contain the intervention or shock definition,
normalization, observation frequency/operator, allowed contemporaneous
structure, latent information, valid moments and sampling regime.
The estimation layer can share lag construction, sparse regressions,
filtering and score reductions, but must preserve which restrictions
identify a structural response.

For many series, sparse VARs, low-rank factors and shared nuisance
learners impose different approximation restrictions. Reducing compute
by screening lag pairs can remove a necessary adjustment path; reducing
frequency can destroy identifiability; pooling countries can impose a
common mechanism that their policy regimes do not share. Scale is thus
a joint scientific and computational design question. Chapter~\ref{r3:discovery} examines what causal discovery can contribute
before those restrictions are used to estimate an effect.

The remaining comparisons follow the response from inference to a banking
calculation. Persistence first changes the reliability of covariance and
weak-instrument procedures. The funding example then carries learned
controls and common calendar shocks into a monetary contrast. Finally,
the dependent-learner example states which restrictions must survive when
that estimator becomes a reusable operation.



% SOURCE docs/release/2026-10-29/source/docs/writing-repair/2026-10-29/src/4a3ed2aed6e7e07f.tex CHAR 0
\section{A practical hierarchy for persistence and dependent inference}
\label{s5:response-inference}
A treasury analyst can reasonably want a deposit response at horizons zero
through twelve, a simultaneous statement about that whole curve, and an
answer robust to a weak monetary-policy instrument. These are three
requirements. Changing the covariance bandwidth addresses dependence;
changing to a reduced-form projection may change the target; neither action
establishes the instrument's exclusion restriction.

For the direct-IV route, stack the raw moments as
\[
 m_t=Z_t\big((X_{i,t+h})_{(i,h)\in\mathcal I\times\mathcal H}',X_{1t}\big)',
 \quad \mu=\E[m_t],\quad \theta_j=\mu_j/\mu_d,
 \quad r_t(\theta)=[I_{d-1},-\theta]m_t.
\]
Here $d-1$ is the number of response coordinates, $X_{1t}$ supplies the
impact normalization, and $j$ denotes an outcome-horizon pair. Relevance,
lead/lag exclusion and the structural linear response argument above make
these ratios causal. A fixed-parameter stationary CLT gives
$\sqrt n(\bar m-\mu)\Rightarrow N(0,\Omega)$.
For each coordinate, invert the residual-moment test
\begin{equation}
 n(\bar m_j-\vartheta\bar m_d)^2\le
 c^2\{\hat\Omega_{jj}-2\vartheta\hat\Omega_{jd}
                +\vartheta^2\hat\Omega_{dd}\}.
 \label{s5:residual-set}
\end{equation}
Using $c=\Phi^{-1}\{1-\alpha/[2(d-1)]\}$ gives a Bonferroni whole-curve
set when every marginal test has the asserted size. No division by
$\bar m_d$ is needed. The answer can be disconnected, unbounded, or empty;
a plotting system must retain that topology. This is a rectangular joint
construction, not an efficient ellipsoidal region or a uniform theorem over
all weakly dependent processes.

\subsection{What prefiltering assumes and what it changes}
The Bartlett long-run covariance estimator uses
$\hat\Gamma_0+\sum_{\ell=1}^L(1-\ell/(L+1))
(\hat\Gamma_\ell+\hat\Gamma_\ell')$ \citep{newey1986}. Its positive
semidefiniteness does not imply a good finite-sample approximation near a
unit root. For centered column scores $u_t=m_t-\mu$, define
$e_t=u_t-Au_{t-1}$. If the covariance sums exist absolutely,
\[
 \Omega_e=(I-A)\Omega_u(I-A)',\qquad
 \Omega_u=(I-A)^{-1}\Omega_e(I-A)^{-T}.
\]
To see this, sum $e_t$ over $t=1,\ldots,n$:
$(I-A)\sum_tu_t+A(u_n-u_0)$. The boundary term is $o_p(\sqrt n)$ in a
stationary finite-variance model. Taking limiting covariances gives the
identity. The residuals need not be white: a misspecified stable AR(1)
prefilter can still be followed by a consistent residual HAC estimator.

Feasible consistency needs more than this identity. For fixed dimension,
assume $\hat A\to A$, $I-A$ is uniformly nonsingular in the chosen fixed
model, and the HAC estimator using estimated residuals differs by $o_p(1)$
from the population-residual estimator. A sufficient replacement condition
under bounded sample quadratic moments is
$L\|\hat A-A\|=o_p(1)$; combine it with the HAC dependence/moment and
bandwidth conditions. Root-$n$ estimation of $A$ and $L=o(\sqrt n)$ give
one route. Continuity of matrix inversion then completes the argument.
An estimated radius below one is a useful numerical check, not evidence
of a uniform stability margin for the true bank process. Near singularity,
recolouring magnifies covariance error by the square of the inverse norm.

The implementation fits a VAR(1) to the score vector, estimates a residual
Bartlett covariance and uses linear solves for recolouring. It reports the
prefilter radius and condition number and refuses an unstable estimate.
The three compared methods use the \emph{same} origins, instruments,
normalization and test inversion. Only covariance construction changes.
This preserves the intervention while revealing the inferential sensitivity.

\begin{center}\small
\begin{tabular}{lrrr}\toprule
External-IV regime & $L=\lfloor n^{1/4}\rfloor$ & $4L$ & AR prefilter\\\midrule
Stationary & .940 & .928 & .935\\
Persistent instrument noise & .905 & .935 & .923\\
Near-unit dynamics & .758 & .933 & .775\\
Near-unit, longer series & .803 & .973 & .820\\
Weak instrument & .910 & .963 & .983\\
Invalid instrument & .240 & .363 & .273\\\bottomrule
\end{tabular}\end{center}
These are 400-repetition whole-curve coverages over 26 coordinates. Weak-set
coverage may improve by becoming unbounded, so coverage must be read with
the archived unbounded-set frequencies and failure diagnostics. In the weak
case, both longer-window and prefiltered constructions are unbounded in at
least one coordinate in 99.25\% of repetitions. The earlier
delta-method curve's 81.8\% near-unit coverage remains an adverse result;
it is not overwritten by this different residual-test comparison. The
longer window is a useful primary challenger for this declared stationary
external-IV family. Prefiltering is not uniformly better. No covariance
method repairs invalid exclusion, as the last row shows.

\subsection{Lag augmentation and a distinct joint construction}
For an AR(1), regress $y_{t+h}$ on $(1,y_t,y_{t-1})$. Partialling out
$y_{t-1}$ makes the effective regressor the innovation in $y_t$. Under the
source's strong innovation restrictions, the relevant score covariance
permits heteroskedasticity-robust inference without a generic residual HAC
penalty. \citet[Definition 1, Assumptions 2--3 and Proposition 1]{lpaug2020}
establishes a marginal result over a specified persistence/root class and
horizon-growth regime. For four prespecified horizons, applying that result
at level $\alpha/4$ and the union bound supplies a simultaneous route;
it does not require independent horizon estimates.

The scalar experiment uses horizons $(1,4,8,12)$, the Eicker--Huber--White
regression covariance and that Bonferroni critical value. Joint 95\%
coverage is .920, .943 and .928 for $\rho=.8,.985,1$ respectively, across
400 repetitions. This is evidence for an AR reduced-form response under
its innovation law. It is not a comparison on the external-IV normalized
target in Equation~\eqref{s5:residual-set}. Uniform marginal validity also
does not by itself permit a growing number of horizons without controlling
the accumulated error \citep{inoue2023}.

The instrument need not affect every date equally. In the inspected 2025
dynamic-LATE preprint, conditional instrument independence, exclusion,
monotonicity, partial first-stage support and smooth response restrictions
make a generalized Wald ratio a weighted marginal effect for the responsive
subpopulation \citep[Assumptions 2.1--2.5 and Proposition 2.1]{casini2025}.
Its further claim of date-specific complier recovery needs deterministic
complier status and local smoothness/comparison assumptions
\citep[Assumptions 2.6--2.8]{casini2025}. Neither statement licenses an
analyst to choose the dates with the strongest observed first stage and use
an ordinary full-sample interval. Selected subsamples change population
weights and require a separately justified selection-aware inferential argument.

\subsection{Choosing a route for a bank's actual clock and panel}
Common shocks require aggregation of customer scores by calendar date;
customer count does not manufacture independent shock realizations. In a
separate AR(.95) date-score mean experiment, short-window HAC covers .500
at 800 dates and .560 at 2,000 dates; prefiltering covers .918 and .958.
This favorable result supports prefiltering for that score model, while
the curve comparison above prevents generalizing it to every response.
With heavy tails, the moment and covariance conditions need checking or a
different robust sampling argument; trimming a shock changes the target
unless a retained-population estimand is specified.

The wider method hierarchy follows the restrictions the bank can defend.
Direct external-IV projections suit a dated external source of variation
with a defensible exclusion. A finite VAR or factor reduction can improve
precision and long-horizon propagation by imposing dynamics and rank; a
comparison with the direct route therefore needs the same shock normalization.
Lag-augmented projections supply a different alternative for their justified
reduced-form or separately identified structural target. An error-correction model is appropriate when levels have
economically justified cointegrating restrictions; estimating rank and
deterministic trends changes the inferential problem and cannot inherit a
stationary score theorem. Differences answer change effects, with level
responses reconstructed by cumulation under the maintained integration
model. Mixed-frequency state-space/MIDAS designs add a measurement and
aggregation law: quarterly output is not three identical monthly observations.
The observation-clock and factor derivations earlier in this chapter explain
the added information and aliasing cost. Breaks require a declared stable
subperiod, a prespecified regime model or explicitly local effects; a
full-sample mean response can otherwise average incompatible policies.

These alternatives are proposed supported model families, not automatically
selected truths. The reusable response operation reports the stability,
normalization, dependence and aggregation premises, the primary estimate,
a substantive challenger and their disagreement. When no family is credible,
it returns an exploratory curve or a sensitivity set with that limitation.
Large-scale reduction remains valuable, but it follows the scientific unit
and target instead of choosing them.


% SOURCE docs/release/2026-10-29/source/docs/writing-repair/2026-10-29/src/50cbe6fc76081094.tex CHAR 0
\section{A bank funding decision with common shocks and learned controls}
\label{s6:funding}
Consider a treasury desk changing the offered deposit rate for one month by
ten basis points. The target is the response of a fixed branch portfolio's
balance stock, followed by the change in funding contribution over seven
months. This is neither a sustained repricing policy nor the response to an
unlabelled reduced-form innovation. Future prices follow their natural rule.
The customer and competitor mechanisms are held fixed; equilibrium repricing
by competitors would require a different model.

For branch $i$, write $B_{it}$ in millions of currency units, $A_t$ in units
of ten basis points, and $W_t$ for an observed common funding driver. A
representative construction is
\begin{align}
 W_t&=\phi W_{t-1}+\sqrt{1-\phi^2}\,\zeta_t,\nonumber\\
 A_t&=a(W_t)+\nu_t,\qquad a(w)=.5w+.2(w^2-1),\nonumber\\
 B_{it}&=\rho B_{i,t-1}+\beta_i A_t+g(W_t)+\eta_t+\epsilon_{it},
 \quad g(w)=1.2w+.4(w^2-1).\label{s6:funding-law}
\end{align}
All innovation sequences are independent Gaussian sequences, with
$\nu_t,\zeta_t$ standardized to variance one; $\eta_t$ is
shared by all branches. Stability margins for $|\phi|,|\rho|<1$, finite
moments and a fixed finite horizon are part of the experiment. The
simulation has 48 equally weighted branches, $\rho=.82$, $\phi=.65$,
$\bar\beta=2$, common noise standard deviation 2 and individual noise
standard deviation 3. Heterogeneous $\beta_i$ range from 1.2 to 2.8.
Adding the baseline stock of 100 per branch changes accounting but not
response estimation.
The simulation starts its centered states at zero and discards 500 dates;
it approximates the stationary law rather than sampling its initial law
exactly. At the declared stability margins the remaining initialization
effect decays geometrically. This construction cannot justify using the
same burn-in for a process arbitrarily close to a unit root.

Let $\bar B_t=\sum_iw_iB_{it}$ for fixed weights summing to one. Because
the intervened price does not change $W$ or future assignment, recursive
substitution gives
\[
 \theta_h=\E[\bar B_{t+h}(A_t+1)-\bar B_{t+h}(A_t)]
       =\bar\beta\rho^h,\qquad h=0,\ldots,H.
\]
This is an additive shift of a continuous action with support on the real
line in the reference law. A bounded bank price grid needs its own feasible
shift and support restriction; Gaussian prices are not a literal product
pricing model. Identification uses the independence of $\nu_t$ from future
structural disturbances conditional on
$X_t=(\bar B_{t-1},W_t)$. That exclusion is stronger than a good prediction
of prices. It fails if officers privately anticipate future withdrawals.
An external instrument is then a substantive alternative only if its own
exclusion and normalization can be defended.

\subsection{The estimator, learning remainder and calendar experiment}
Define $m_h(X)=\E[\bar B_{t+h}\mid X]$ and $e(X)=\E[A_t\mid X]$.
For $v_t=A_t-e(X_t)$, the structural law implies
\[
 \E[v_t\{\bar B_{t+h}-m_h(X_t)-\theta_hv_t\}]=0,
 \qquad q=\E v_t^2=1.
\]
The finite basis $(1,\bar B_{t-1},W_t,W_t^2-1)$ contains these nuisance
functions under Equation~\eqref{s6:funding-law}: Gaussian AR conditional
means of future linear and quadratic driver terms stay in this basis.
The partialling-out construction is a selected specialization of the
orthogonal approach in \citet{dml2018}; the time-series extension in
\citet{ciganovic2026} requires its additional stability and sampling
conditions, not merely chronological input rows.

An independent pilot realization chooses a ridge penalty from
$\{0,1,10,100\}$ by a chronological half-pilot validation loss, then fits
all nuisance functions on the pilot. The family, candidate penalties and
normalization are fixed before that pilot. The evaluation sample never
selects a penalty, horizon, regime or bandwidth. Independent realization
means an independent statistical population under the same law; an adjacent
historical block is not literally independent under this model. In a bank,
using a separated block instead requires a gap/mixing argument and its
finite-sample sensitivity. Pooling another country additionally requires
transport of the nuisance law. The reference deliberately makes this
information cost visible.

With frozen pilot functions $\hat e,\hat m_h$, set
$\hat v=A-\hat e$ and estimate
\[
 \hat\theta_h=
 \frac{\mathbb P_T[\hat v(\bar B_{t+h}-\hat m_h)]}
      {\mathbb P_T\hat v^2}.
\]
Write $d_e=\hat e-e$, $d_h=\hat m_h-m_h$ and
$u_h=\bar B_{t+h}-m_h-\theta_hv$. Conditional on the pilot, the
expectation of the estimating score at $\theta_h$ is exactly
\begin{equation}
 \E[d_e(X)d_h(X)]-\theta_h\E[d_e(X)^2].
 \label{s6:pilot-remainder}
\end{equation}
To see this, expand
$(v-d_e)(u_h-d_h+\theta_hd_e)$; conditional means of $v$ and $u_h$
vanish. Thus independence alone does not remove learning bias. Sufficient
conditions are nuisance consistency and
$\|d_e\|_2(\|d_h\|_2+|\theta_h|\|d_e\|_2)=o_p(T^{-1/2})$,
with corresponding moment/stochastic-equicontinuity bounds. The fixed
correct basis, independent pilot size proportional to $T$ and bounded
penalty candidates give root-pilot coefficient error; the remainder is
$O_p(T^{-1})$. An arbitrary forest or representation learner would need
a new rate argument.

Under the stable Gaussian-driven polynomial law, covariances and the
required mixing/moment sums decay geometrically. The vector influence is
\[
 \psi_{th}=q^{-1}v_tu_{th},\qquad
 \sqrt T(\hat\theta-\theta)
 =T^{-1/2}\sum_t\psi_t+o_p(1).
\]
A mixing central limit theorem therefore gives covariance
$\Omega=\sum_{k\in\mathbb Z}\E[\psi_t\psi_{t-k}']$.
The evaluation residuals replace both $\theta$ and the pilot nuisances in
$\hat\psi$. Bartlett covariance with
$L=\lfloor2T^{1/5}\rfloor$ uses the denominator $T$ at every lag and
estimates $\Omega/T$ \citep{newey1986}. For the fixed basis and moments
here, coefficient replacement contributes at most
$O_p(LT^{-1/2})$ to the lag-sum bound, and the taper and tail biases vanish.
The exponent is strictly below the general source's $1/4$ bound.
This argument is pointwise in the stability margin and fixed $H$; it gives
no uniform near-unit-root guarantee. Overlapping horizons share future
outcomes, and their covariance must remain in $\hat\Omega$.

Increasing branches removes part of $\epsilon_{it}$, but not $\eta_t$ or
the common action innovation. Calendar origins are the information count.
Dividing the calendar variance by the number of branches is an adverse
comparison, not another interpretation of cluster robustness. Bonferroni
normal intervals give a conservative joint response band; a prespecified
financial contrast uses its own scalar normal interval. Simulated coverage
checks the chosen regime, not the asymptotic assumptions.

\subsection{Funding accounting and method decisions}
Let $N$ be the branch count, $c=10^6$ the currency conversion, $m=.002$
the monthly net alternative-funding margin and $d=.99$ the discount factor.
The first-month additional rate is $r=.001/12$. Relative to no shift,
\begin{equation}
 V=Nc\left\{\sum_{h=0}^H d^h m\theta_h
                 -r(B_0+\theta_0)\right\},\qquad
 \widehat{\Var}(\hat V)=k'\hat\Omega k/T,
 \label{s6:funding-account}
\end{equation}
where $k_h=Nc(d^hm-r\mathbf1\{h=0\})$ and $B_0=100$ is fixed.
The higher rate also applies to induced first-month balances. Omitting that
term overstates value. There is no cost saving from relabelling a balance
stock as a transaction flow. If baseline exposure, margin or discounting is
estimated, their influence and covariance must be added; the displayed
interval conditions on these economic inputs.

The primary procedure is the pilot-trained calendar score because its
functional, learning error and financial covariance can all be derived for
this information regime. Direct residual regression with the same frozen
pilot and Bartlett rule is a computational comparator: it solves precisely
the same score. An outcome regression fitted jointly on evaluation data can
be more efficient in the correctly specified finite basis, but is a
different fitted experiment. External-IV projections become the primary
alternative when the price innovation is privately confounded and a valid
instrument is available; their weak normalization needs the set-valued
inference already derived. Lag-augmented innovation projections do not by
themselves solve that IV problem.

A short bank panel with endogenous lagged balances calls for the separate
initial-law and lag-moment analysis, not multiplication of the calendar
sample size. A mixed-frequency analysis can retain daily prices and monthly
stock measurements, but must specify whether the monthly stock sees the
last daily action, its path or an integral. A MIDAS distributed-lag
restriction reduces coefficients only after that observation equation is
fixed. Synthetic controls address a sustained unit intervention under donor
restrictions; they do not automatically recover this one-month innovation
response. Each is a substantive alternative when its observation and
assignment premises match the bank's question.

If evaluation outcomes suggest a break, choosing the favorable subperiod and
reusing this interval is invalid. The result must become exploratory, or
the bank must define a new regime-specific target and obtain independent
evaluation information. The public operation binds the pilot, calendar,
bandwidth, observation mapping, price normalization and financial weights
to its inference identity. This is a necessary consequence of the derivation,
not merely a software cache preference.


% SOURCE docs/release/2026-10-29/source/docs/writing-repair/2026-10-29/src/63f4d4a9da080ca4.tex CHAR 0
\section{A dependent learner with a public scientific contract}
\label{c7:blocked-operation}
The separated-block construction of Section~\ref{s4:dependent-learning}
can support an operation only if its learning and calendar assumptions survive
the interface. The target remains the partially linear IV coefficient with
score $(Z-m(C))\{Y-\ell(C)-\theta(D-r(C))\}$. Its population slope must be
nonzero; controls and the instrument must satisfy the substantive restrictions
before any nuisance prediction. The product does not infer those restrictions
from a high validation score.

The request names a bounded geometrically absolutely regular calendar law,
a lower bound $c$ on the exponential decay rate, feature-lag and outcome-horizon
spans, a separation gap, a nuisance-rate statement and one of the implemented
learners. The two evaluation blocks are separated after expanding their full
lag/outcome windows. The operational check requires
$g\ge\lceil3\log n/c\rceil+H+L_C$. The supplied $c$ is a population
assumption; an empirical autocorrelation plot does not establish it. Arbitrary
user estimator objects are not accepted as qualified learners.

Each block is evaluated using nuisance fits trained on the other block. Lasso
tuning uses only a further separated training/validation split inside that
training block. Three prespecified penalties share the order
$\sqrt{\log n\log(np)/n}$; predictions are clipped to the declared known
envelope. The fixed-dictionary challenger uses ordinary least squares on the
same controls. The fitted score retains the exact product remainder derived
earlier, and the covariance still requires $L_n\delta_n=o(1)$ in addition to
the orthogonality product condition. The operation uses $L_n=\lfloor n^{.2}\rfloor$.
With a fixed correct sparse dictionary these conditions follow from the
earlier sufficient rate construction; accepting a learner name alone does
not establish approximation error or restricted curvature. The general score
conditions of \citet[Assumptions 3.1--3.2]{dml2018} remain conditions, not an
automatic dependent-data theorem.

Scores use one global centering mean. Their Bartlett lag products are summed
within each evaluation block with the total retained sample as denominator;
the deleted gap is not collapsed into adjacency. The weak-IV extension is a
different inference operation. A numerically nonzero sample slope does not
prove the fixed nonzero population slope required by this regular method.

The executed example uses 3,000 dates, 20 bounded controls and a 90-date gap.
Across 200 independent simulations, nominal 95\% coverage is $187/200=.935$
for the sparse learner, $.930$ for the fixed basis and $.940$ for an oracle
with the true conditional means. The sparse coverage Monte Carlo interval is
$[.8914,.9649]$; the experiment is not precise evidence for exact calibration.
The intercept-only challenger covers in none of the 200 draws. When both
outcome and instrument nuisance functions omit a common nonlinear interaction,
both fitted learners again cover in none of the draws, while the oracle covers
in $187/200$. This is the consequence of a persistent product remainder, not
an error corrected by adding a HAC covariance after a misspecified fit.

A bank can therefore reuse this selected learning route with an explicit
dictionary, causal information and rate argument. It cannot substitute a
different learner on the strength of a common method signature. Nor can the
independent-pair audit result be attached to this calendar score without a
joint sampling derivation.

